%----------
%	OBJECTIVES AND MOTIVATION
%----------	

Esta sección está dedicada a la presentación de los objetivos generales y específicos que se pretenden lograr con el desarrollo de este Trabajo Final de Máster. Podemos identificar los siguientes objetivos:

\textbf{OG1.} Investigar las arquitecturas y métodos de diseño, entrenamiento y refinado de los modelos de IA más avanzados, incluídos los LLM y sistemas basados en agentes, para analizar su viabilidad y prestaciones para detectar y localizar variables cuantitativas y cualitativas relacionadas con la representación mediática de las mujeres y susceptibles de ser integradas posteriormente en una herramienta que apoye la postedición de noticias.

\textbf{OE1.} Diseñar la arquitectura conceptual del sistema, definiendo los roles funcionales de los agentes, su relación con los flujos de trabajo periodísticos y los mecanismos de interacción entre los distintos módulos (análisis, verificación y recomendación).

\textbf{OE2.} Desarrollar y evaluar estrategias de prompting y herramientas asociadas (criterios, plantillas, contextos expertos) para la identificación coherente de variables discursivas, estilísticas y de sesgo.

\textbf{OE3.} Implementar técnicas de representación y recuperación de información, incluyendo RAG, búsqueda semántica y medidas de similitud, para mejorar el razonamiento contextual de los modelos y enriquecer sus recomendaciones.

\textbf{OE4.} Comparar distintos modelos de lenguaje y pipelines de procesamiento del lenguaje natural a fin de determinar su capacidad para detectar sesgos explícitos e implícitos y su idoneidad para su integración en un sistema de postedición en el futuro.

\textbf{OE5.} Proponer directrices para la futura integración del sistema en entornos editoriales reales, incluyendo escenarios de uso, requisitos éticos, criterios de explicabilidad y mecanismos de supervisión humana.


\newpage